% 1. Preamble
\documentclass[12pt, a4paper]{article}

% Import Packages
\usepackage[utf8]{inputenc}        % Encoding
\usepackage{amsmath, amssymb}     % Advanced math formulas
\usepackage{mathtools}           % Math enhancements
\usepackage{graphicx}            % Image insertion
\usepackage{hyperref}            % Interactive hyperlinks & PDF bookmarks
\usepackage{color}               % Color text and backgrounds
\usepackage{geometry}            % Page geometry
\usepackage{comment}             % Multi-line comments
\usepackage{titling}             % Custom title formatting

\geometry{a4paper, margin=1in}

% Custom Command Definitions
\newcommand{\Sball}{\mathcal{S}}
\providecommand{\norm}[1]{\left\lVert#1\right\rVert}

% Hyperref Setup for Clean PDF Links
\hypersetup{
    colorlinks=true,
    linkcolor=blue,
    urlcolor=blue,
    pdftitle={Pattern Arithmetic},
    pdfauthor={Sandeep Lakshminarayan Chiluveru}
}

% Metadata & Custom Title Banner
\pretitle{\begin{center}\Huge\fontfamily{tt}\selectfont ================================================================================\\[0.5em]}
\posttitle{\\[0.5em] ================================================================================\end{center}}

\title{Pattern Arithmetic: Phase-Colored Geometric Framework for High-Density Information via Unit Probability Spheres}
\author{Sandeep Lakshminarayan Chiluveru}
\date{\today}

\begin{document}

\maketitle

% =====================================================================
% TODO -- MASTER PUNCH LIST (review pass, keep until resolved)
% =====================================================================
% [T1] State vs. field conflation (Axiom 1 vs. everywhere else).
%      Axiom 1 defines S as a probability FIELD (triple integral over
%      the whole sphere = 1). Axioms 4-6, T, and every worked example
%      treat a pattern as a single POINT (r,theta,phi,chi). Pick one
%      object and rewrite the other axioms/definitions consistently,
%      or explicitly define the map between "a point" and "a field"
%      (e.g. a point is a delta-function field) and use it.
% [T2] Axiom 2's quantifier "for all (theta,phi,chi)" applied to a
%      single vector's norm is ill-formed -- a single vector doesn't
%      range over theta,phi,chi independently. Rewrite as a plain
%      condition on r (or on the field, once T1 is settled).
% [T3] Axiom 3 (S = infinity, white) contradicts the paper's own
%      framing of a BOUNDED register with capped r and idempotent
%      ⊕. Under idempotence + capping, infinite union of distinct
%      states should saturate (e.g. r -> 1, full coverage), not
%      diverge to infinity. Either drop this axiom, or replace it
%      with a saturation statement, or explain precisely what norm
%      is being taken to infinity and why boundedness doesn't apply.
% [T4] Axiom 5 (interference) gives ONLY a magnitude formula. It does
%      not say how the resultant direction (theta,phi) or phase chi
%      is computed when A and B point in different directions. The
%      "blended phase = 60 deg" example in the Introduction and the
%      "chi = 66.6 deg" example in Sec. 5.1 are NOT reproducible from
%      any formula actually stated in this paper (checked: simple
%      average, magnitude-weighted average, and vector/phasor sum all
%      give different numbers -- 60, 68.6/72, and 73.9 respectively,
%      none of which is 66.6). Either derive the phase/direction
%      update rule explicitly and re-derive the examples from it, or
%      relabel these as illustrative/unverified placeholders.
% [T5] The offset operator (Sec. 4) is only demonstrated on a 1-D
%      great circle (man/woman/king/queen all share one angular
%      coordinate). Subtraction/addition of points is not a
%      well-defined operation on a curved 2-sphere in general -- it
%      needs a stated construction (e.g. lift to tangent space via
%      log map, transport, then exp map) before the "offset" claim
%      extends beyond this degenerate 1-D case.
% [T6] Sections 6-8 (Operational Algebra numerics beyond 5.1,
%      Real-World Models, Discussion & Related Work) are mostly
%      placeholder headers/bullets not derived from the axioms above.
%      Flesh out or clearly mark as future work; see inline TODOs.
% [T7] Missing related-work comparisons: RotatE (relations as complex
%      rotations -- close to the T operator), Bloch sphere qubit
%      representation (already named as a header, undiscussed), and
%      von Mises-Fisher / hyperspherical distributions for directional
%      probability data. The Discussion section currently has no
%      content under any of its three headers.
% =====================================================================

% =====================================================================
% GAPS 
% =====================================================================
\begin{comment}
Gaps, in order
A. Close $\oplus$: $A\oplus B=(r',\theta',\phi',\chi')$, with a rule when directions differ; one numerical check.
B. Rewrite Axiom 3 (White = unit full-spectrum, not $\infty$). Kill the $66.6^\circ$ line.
C. One sentence: the atom is a 4-tuple or a field.
D. Offset on the manifold: $\hat n_q \propto \hat n_k-\hat n_m+\hat n_w$, then say what $r$ and $\chi$ do.
E. One use case through the official four steps, with numbers. Storm first.
F. One numeric encoding; show $7+5=12$ and that $\oplus$ on the same marks is not $12$.
G. Related work as comparison paragraphs (Bloch, word2vec/cosine, Riemann sphere, $\mathrm{SO}(3)$, spherical harmonics).
H. Leave the concept dictionary in Future Work until A–F exist.
\end{comment}
% =====================================================================


% ABSTRACT
\begin{abstract}
Ordinary arithmetic treats the integer 1 as a unit of truth: $1=1$ and $1+1=2$. Counting is only one algebra. Colors and musical notes are also units of truth; combining them yields a hue or a chord, not a larger integer. This paper treats a pattern as a higher-order unit and represents one pattern by one state on a unit probability sphere $\mathcal{S} \subset \mathbb{R}^3$, with magnitude $r \in [0, 1]$, direction $(\theta, \phi)$, and internal phase-color $\chi \in [0, 2\pi)$. Pattern Arithmetic is the algebra of such units: identical states do not double after normalization ($A \oplus A = A$), and opposite phases cancel ($A - A = 0$). A second operator $T$ acts by rotation and phase shift, allowing the same ingredients to form different ordered expressions. Words and numbers occupy the same sphere under different encodings and meet in a document by order and typed action, not by blending. The aim is a single bounded register that carries overlapping identities, structural paths, and mixed prose without a discrete grid.
\end{abstract}

\section{Introduction}

Ordinary arithmetic uses the integer $1$ as a unit of truth.
The identity $1=1$ says that two tokens are of the same kind.
The rule $1+1=2$ says that combining two such tokens produces a
larger quantitative state. That algebra answers the question
\emph{how many?}

It is not the only algebra. A color and a musical note are also
units of truth, but their combination does not enlarge a pile.
Under additive light, red and green yield yellow. Two copies of
the same yellow, once brightness is normalized, remain yellow.
Two notes yield a chord, not a larger integer; a tone plus its
opposite phase yields silence. In each case the word ``add''
names a different rule, and therefore a different kind of unit.

When the object of interest is a structured whole rather than a
count, a hue, or a pitch, the natural unit is a \emph{pattern}.
This paper represents one pattern by one state on a unit
probability sphere $\mathcal{S} \subset \mathbb{R}^3$:
\begin{itemize}
    \item $r \in [0, 1]$ is the strength of the unit, capped so that a state remains a single unit of truth;
    \item $(\theta, \phi)$ is the direction of the pattern on the sphere;
    \item $\chi \in [0, 2\pi)$ is an internal phase-color, so two identities may occupy the same direction without overwriting one another.
\end{itemize}
Pattern Arithmetic is the algebra of these units. Identical states do not double after normalization ($A \oplus A = A$). Opposite phases cancel ($A - A = 0$). Rotations of the sphere preserve the unit.

A single numerical picture is enough to see the difference from counting. Take two unit states at the same place, one with phase $0^\circ$ (red) and one with phase $120^\circ$ (green). Counting would return $2$. Phase-aware combination returns a single unit whose blended phase is $60^\circ$ (yellow).
% TODO [T4] No rule for computing the resultant phase/direction has
% been given yet at this point in the paper (Axiom 5, introduced
% later, only yields a magnitude). Either forward-reference the exact
% equation that produces 60 deg here, or state plainly that this is
% an illustrative target the later formal machinery does not yet
% reproduce in the general (unequal-strength, off-axis) case.
The same two strengths with opposite phase ($0^\circ$ and $180^\circ$) return the null state, black. Two units of truth become either one new identity or none.

The working order of the paper is not that mix alone. It is a three-step spine:
\[
\text{Collapse (bowl)} \;\longrightarrow\; T\text{ (path)} \;\longrightarrow\; \text{Result}.
\]
Collecting red and green is not yet yellow, and not yet a poem.
Yellow is what you get if the later path is $T_{\mathrm{paint}}$.
Red and green may equally well walk a $T_{\mathrm{poem}}$ path; the
result need not be yellow --- though a given poem may still arrive
there. The pathway, not the shells, chooses the result. A verse is
another path.
Each word of ``Twinkle, twinkle, little star\ldots'' is a unit state
on the same sphere. Collapse yields the bag of distinct words
($\texttt{twinkle}$ appears twice in the line and only once in the
bowl). The poem is the path that walks those states in order,
including the repeated $\texttt{twinkle}$. The bag is not the song.
Numericity is not stored by $\oplus$. Two identical units collapse
to one address; that is the point of $A\oplus A=A$. How many times
the address is \emph{visited} lives in the path (the tape says
\texttt{twinkle} twice). How many as a \emph{count} lives in
$\mathtt{num}$ mode, where $1+1+1+1=4$ and four tokens of $1$ are
not collapsed. Strength $r$ is a poor substitute for either: a
second \texttt{twinkle} cannot honestly raise $r$ past $1$, and
``four ones'' is schoolbook addition, not a fuller unit of truth.
The same bowl may be read more than one way. Loading is not meaning.

\subsection{The Limitations of Discretized Matrix Grids}
A discrete matrix grid is the wrong box for such a unit. Traditional matrix representations of continuous state spaces suffer from two severe structural bottlenecks:
\begin{itemize}
    \item \textbf{Memory Scaling Bottlenecks:} A voxel cube of resolution $N$ stores $O(N^3)$ cells, growing steeply with resolution.
    \item \textbf{Signal Identity Destruction:} Mapping multiple overlapping variables to identical spatial coordinates often results in the destructive overwriting of underlying signal identities.
\end{itemize}

\subsection{The Phase-Colored Unit Probability Sphere Framework}
To resolve these grid constraints, the sphere is offered as one bounded object that can carry overlapping identities by direction and phase, without enlarging a lattice.
\begin{itemize}
    \item \textbf{Spatial Orientation $(\theta, \phi)$:} Defines continuous directional coordinates across the spherical surface.
    \item \textbf{Magnitude Likelihood ($0 \leq r \leq 1$):} Represents bounded probability density within the sphere.
    \item \textbf{Internal Color Phase ($\chi \in [0, 2\pi)$):} Serves as an independent phase degree of freedom.
\end{itemize}

\subsection{Core Contributions \& High-Level Intuition}
The principal contributions of Pattern Arithmetic are:
\begin{itemize}
    \item \textbf{High-Density Superposition:} Storing overlapping state vectors concurrently within a single bounded spatial manifold.
    \item \textbf{Parameter-Efficient Field Storage:} Eliminating discrete voxel bounds by treating states as continuous probability distributions.
\end{itemize}

Later sections distinguish a numeric encoding from a lexical encoding on the same sphere, so that a report may carry words and numbers together without forcing them through $\oplus$.

The next section turns these definitions into formal mathematical axioms.


\section{Axiomatic Foundations Of Pattern Arithmetic}

\subsection{Process principle: Collapse, path, result}
\label{subsec:process_principle}
Pattern Arithmetic proceeds in three named acts, kept separate so that a bowl may grow before it is read:
\begin{equation}
    \underbrace{A_1 \oplus A_2 \oplus \cdots \oplus A_n}_{R\ =\ \text{the bowl}}
    \;\longrightarrow\;
    T_{w_k}\circ \cdots \circ T_{w_1}(R)
    \;\longrightarrow\;
    \text{Result}.
    \label{eq:pipeline}
\end{equation}
Here $\oplus$ \emph{collects}: identical tokens do not double ($A \oplus A = A$), and opposite marks may sit together until a later reading. $T$ is an ordered path, not a mix. The result is whatever that path (or another named reading) makes of $R$ --- a poem, a report, a staged endpoint.

Paint is one species of path, written $T_{\mathrm{paint}}$:
\[
T_{\mathrm{paint}}(\{\mathrm{red},\mathrm{green}\}) = \mathrm{yellow}.
\]
The same pair may take $T_{\mathrm{poem}}$ instead. That result is a
verse about the two colors, not a hue --- unless the poem itself
walks back to yellow. Shells do not own a pathway. A verse is another
species. One bowl, several pathways. One may collect
and read in a single gesture; the algebra does not force the merge.
This principle is an axiom of \emph{order of acts}, not a new formula
for a 4-tuple.

\subsection{Axiom 1: The Normalized Unit Truth State ($\mathcal{S} = 1$)}
\label{subsec:axiom1}
To ensure stability across multi-variable operations, every state vector field $\vec{v} = (r, \theta, \phi, \chi)$ on the Unit Probability Sphere $\mathcal{S} \subset \mathbb{R}^3$ represents an immutable unit of truth ($\mathcal{S} = 1$). Global likelihood integrates strictly to unity:
\begin{equation}
    \iiint_{\mathcal{S}} P(r, \theta, \phi) \, dV = 1.0
    \label{eq:unit_truth}
\end{equation}
% TODO [T1] This axiom defines a state as a PROBABILITY FIELD
% P(r,theta,phi) integrating to 1 over the whole sphere. Every other
% axiom, the T operator, and every worked example instead treat a
% state as a single POINT (r,theta,phi,chi). These are different
% mathematical objects. Decide which one "a pattern" actually is, and
% either (a) rewrite this axiom as a normalization on the 4-tuple
% (e.g. r in [0,1] plus a statement about what varies), or (b) keep
% the field definition and explicitly define points as a special case
% (e.g. delta-function fields), then use that mapping consistently
% everywhere a "state" is combined, rotated, or compared.

\subsection{Axiom 2: The Absolute Null State ($\mathcal{S} = 0$, Black)}
\label{subsec:axiom_null}
The state $\mathcal{S} = 0$ defines the absolute lower bound of the state space, corresponding to complete destructive interference, total phase cancellation, or zero-state equilibrium:
\begin{equation}
    \norm{\vec{v}} = 0 \quad \forall \; (\theta, \phi, \chi) \implies \mathcal{S} = 0 \quad \text{(Black)}
    \label{eq:null_state}
\end{equation}
% TODO [T2] "for all (theta,phi,chi)" is not a meaningful condition on
% a single vector v -- a given vector doesn't range over theta,phi,chi
% independently; its norm is just norm(v), full stop. Either (a) state
% this simply as "r = 0" (direction and phase are then undefined,
% which is itself worth a sentence), or (b) if this is meant as a
% condition on the FIELD P(r,theta,phi) from Axiom 1 (P = 0
% everywhere), say that explicitly and make Axiom 1 vs. point-state
% consistent first (see T1).
Physically and visually, $\mathcal{S} = 0$ represents an empty probability field devoid of active signal energy or directional orientation.

\subsection{Axiom 3: The Infinite Superposition State ($\mathcal{S} = \infty$, White)}
\label{subsec:axiom_infinite}
The state $\mathcal{S} = \infty$ defines the theoretical upper bound of maximal entropy, representing total spectral superposition across all continuous orientations $(\theta, \phi)$ and internal color-phase frequencies $\chi \in [0, 2\pi)$:
\begin{equation}
    \mathcal{S} = \lim_{N \to \infty} \bigoplus_{i=1}^{N} A_i \implies \mathcal{S} = \infty \quad \text{(White)}
    \label{eq:infinite_state}
\end{equation}
% TODO [T3] This directly conflicts with the paper's central framing
% (abstract, Sec. 1.2) of S as a BOUNDED register: r is capped at 1,
% and Axiom 4 makes ⊕ idempotent. Idempotent + capped composition of
% infinitely many distinct states should SATURATE (e.g. full angular
% coverage at r -> 1), not diverge to an unbounded S = infinity. As
% written, no norm is even defined for "infinity" on a unit ball, so
% the equation isn't dimensionally meaningful yet. Either (a) replace
% this axiom with a saturation statement consistent with idempotence,
% or (b) if "white" is meant purely as a qualitative pole (like black
% in Axiom 2) rather than a literal state reachable via ⊕, say so and
% drop the limit equation, which currently implies something the rest
% of the algebra can't actually produce.
Physically and visually, $\mathcal{S} = \infty$ corresponds to full-spectrum white light, where all possible state vectors co-exist simultaneously without spatial exclusion.

\subsection{Axiom 4: Information Idempotence ($1 \oplus 1 = 1$)}
\label{subsec:axiom_idempotence}
Combining identical information states yields no expansion of likelihood volume. Upon intensity normalization:
\begin{equation}
    A \oplus A = \frac{A + A}{\norm{A + A}} = A
    \label{eq:idempotence_formal}
\end{equation}

\subsection{Axiom 5: Phase-Aware Superposition \& Interference}
\label{subsec:axiom_interference}
This rule is a \emph{reading} of a collected bowl (the Mix of the process principle), not the act of loading. The combination of two truth states $A$ and $B$ already in the bowl is governed by relative phase difference $\Delta\chi$:
\begin{equation}
    A \pm B = \sqrt{\norm{A}^2 + \norm{B}^2 + 2\norm{A}\norm{B}\cos(\Delta\chi)}
    \label{eq:interference_formal}
\end{equation}
% TODO [T4] This equation gives only the MAGNITUDE of the combined
% state. It says nothing about the resultant direction (theta,phi)
% when A and B point different ways on the sphere, or the resultant
% phase chi in the general (unequal-magnitude) case. Without that,
% the "blended phase = 60 deg" example in Sec. 1 and the
% "chi = 66.6 deg" example in Sec. 5.1 have no derivation -- I checked
% several natural candidates (simple angular average, magnitude-
% weighted average, phasor/vector-sum angle) against the 66.6 deg
% figure and none reproduce it. Add the missing direction/phase
% update rule (e.g. weighted circular mean, or phasor addition treating
% (r,chi) as a complex amplitude, made explicit) and re-derive both
% worked examples from it, or mark those numbers as illustrative
% targets rather than computed results.
Destructive anti-phase interference ($\Delta\chi = 180^\circ$) collapses composite states directly into the Null State ($A - A = 0$, Black) defined in Axiom~\ref{subsec:axiom_null}.

\subsection{Axiom 6: Rotational Conservation (\texorpdfstring{$\times \mathbf{R}$}{x R})}
\label{subsec:axiom_rotation}
Multiplication operates as a Lie group rotation $\mathbf{R} \in \mathrm{SO}(3)$ in continuous spherical space, preserving total truth volume:
\begin{equation}
    \norm{\mathbf{R} \cdot A} = \norm{A} = 1.0
\end{equation}

\section{Transformation Algebra: From Identity to Expression}
\label{sec:transformation_algebra}

\subsection{Motivation: Two Layers of Meaning}
\label{subsec:transform_motivation}
The operators established so far ($\oplus$, $\pm$) answer the question
\emph{what is this state?} Combining Sky, Sun, and Wind under $\oplus$
yields a single, order-independent resultant --- a new composite unit
of truth, a bag of ingredients. This is a \emph{lexical} operation: it mints blended identities, the way
a palette mints a hue from pigments. Analogies such as ruler-man-king are \emph{not} blends; they are deferred to the offset operator of Section~\ref{sec:analogical_offset}. Blend $\oplus$ does not, and structurally
cannot, distinguish a poem in which the sun rises through a wind-blown
sky from a poem in which wind extinguishes a darkening sun --- both
share the identical resultant
$R = \text{Sky} \oplus \text{Sun} \oplus \text{Wind}$, since $\oplus$
is commutative and idempotent by Axiom~\ref{subsec:axiom_idempotence}.

The distinction between these two poems is not a difference in
\emph{ingredients}; it is a difference in \emph{how one state acts
upon another, and in what order}. This is a \emph{relational}
operation, and it requires a second operator, $T$, layered on top of
$\oplus$: where $\oplus$ builds nouns, $T$ builds sentences. A poem is
not the resultant $R$ itself, but a specific trajectory traced across
$\mathcal{S}$ by successive transformations applied to $R$.

\subsection{Formal Definition of the Transformation Operator}
\label{subsec:transform_definition}
Rather than introduce $T$ as an unrelated new primitive, it is defined
using the same rotational machinery already given in
Axiom~\ref{subsec:axiom_rotation}: every state $w \in \mathcal{S}$
acts on another state $x \in \mathcal{S}$ by rotating $x$'s direction
about $w$'s own axis, by an angle proportional to $w$'s strength, while
additively shifting $x$'s phase by $w$'s phase:
\begin{equation}
    T_w(x) \;=\; \Big(r_x,\;\; \mathrm{Rot}_{\hat n(\theta_w,\phi_w)}\!\big(\alpha(r_w)\big)\cdot(\theta_x,\phi_x),\;\; (\chi_x + \chi_w) \bmod 2\pi\Big)
    \label{eq:transform_def}
\end{equation}
where $\hat n(\theta_w,\phi_w)$ is the unit axis defined by $w$'s
direction, $\mathrm{Rot}_{\hat n}(\alpha) \in \mathrm{SO}(3)$ is the
rotation of Axiom~\ref{subsec:axiom_rotation} about that axis by angle
$\alpha$, and $\alpha(r_w) = \pi r_w$ is the simplest monotonic choice
mapping strength to rotation angle (a full-strength unit, $r_w = 1$,
rotates by the maximal $180^\circ$). Magnitude $r_x$ is left invariant:
a transformation reorients and recolors a state, it does not
strengthen or weaken it --- that remains the role of $\oplus$.

\subsection{Derived Properties}
\label{subsec:transform_properties}
The following properties follow from Equation~\eqref{eq:transform_def}
rather than being separately postulated.

\paragraph{Identity.} If $w = \mathbf{0}$ (the Null State of
Axiom~\ref{subsec:axiom_null}), then $\alpha(0) = 0$ and the phase
shift is $0$, so $T_{\mathbf{0}} = \mathrm{id}$.

\paragraph{Invertibility.} Define the inverse state $w^{*} = (r_w, \pi - \theta_w, \phi_w + \pi, -\chi_w \pmod{2\pi})$. 
Then $T_{w^{*}}\!\big(T_w(x)\big) = x$ for all $x$. Because the magnitude $r_w$ remains strictly positive, reversing the transformation requires rotating by the same angle $\alpha$ around the antipodal axis $-\hat{n}(\pi - \theta_w, \phi_w + \pi)$, while the inverse phase cancels the internal color shift. Every transformation has a well-defined inverse. Note that this is distinct from anti-phase interference ($A - A = 0$, Axiom~\ref{subsec:axiom_null}), which annihilates a state; $T_{w^{*}}$ instead \emph{undoes a transformation} on a state that persists.

\paragraph{Associativity.} Since each $T_w$ is a function
$\mathcal{S} \to \mathcal{S}$, composition of transformations is
associative automatically, by the associativity of function
composition; no additional axiom is required.

\paragraph{Non-commutativity.} Rotations about distinct axes in
$\mathrm{SO}(3)$ generically do not commute: for $a, b$ with distinct
directions $(\theta,\phi)$,
\begin{equation}
    T_a\big(T_b(x)\big) \;\neq\; T_b\big(T_a(x)\big).
    \label{eq:noncommute}
\end{equation}
This is not an additional assumption; it is inherited directly from
the well-established non-commutativity of $\mathrm{SO}(3)$, already
invoked by Axiom~\ref{subsec:axiom_rotation}. It is this property,
and only this property, that allows two poems built from the same
resultant $R$ to diverge: order of application is meaningful even
though the underlying identity is not.

\subsection{Poems as Transformation Trajectories}
\label{subsec:transform_trajectories}
A poem $\pi$ built from resultant $R$ and an ordered sequence of
transforming states $w_1, w_2, \ldots, w_n$ is the trajectory
\begin{equation}
    \pi = \big(R;\, w_1, \ldots, w_n\big): \quad
    R \;\longrightarrow\; T_{w_1}(R) \;\longrightarrow\; T_{w_2}T_{w_1}(R) \;\longrightarrow\; \cdots \;\longrightarrow\; T_{w_n}\cdots T_{w_1}(R).
    \label{eq:poem_trajectory}
\end{equation}
Two poems may share an identical resultant $R$ (the same theme, the
same ``ingredients'') while tracing entirely different trajectories
--- this is the formal sense in which many poems arise from the same
units of truth. Conversely, because rotation composition is
many-to-one, two distinct trajectories may terminate at the same
endpoint; a poem's identity is therefore carried by its full path
through $\mathcal{S}$, not by its final state alone.

\subsection{Worked Numerical Example: Sun, Wind, and Sky}
\label{subsec:transform_example}
Let Sky be the state acted upon, with direction
$(\theta,\phi) = (45^\circ, 45^\circ)$, i.e.\ Cartesian axis
$(0.500, 0.500, 0.707)$. Let Sun define an axis along
$(\theta,\phi) = (90^\circ, 0^\circ)$, i.e.\ $(1,0,0)$, with strength
$r_{\text{Sun}} = 1.0$, giving rotation angle $\alpha_{\text{Sun}} =
180^\circ$. Let Wind define an axis along $(\theta,\phi) = (90^\circ,
90^\circ)$, i.e.\ $(0,1,0)$, with strength $r_{\text{Wind}} = 0.7$,
giving $\alpha_{\text{Wind}} = 126^\circ$.

\emph{Order 1 --- Sun acts first, then Wind} (the sun rising through a
wind-blown sky):
\begin{equation}
    T_{\text{Wind}}\big(T_{\text{Sun}}(\text{Sky})\big) \;\approx\; (-0.866,\, -0.500,\, 0.011)
\end{equation}

\emph{Order 2 --- Wind acts first, then Sun} (wind extinguishing a
darkening sun):
\begin{equation}
    T_{\text{Sun}}\big(T_{\text{Wind}}(\text{Sky})\big) \;\approx\; (0.278,\, -0.500,\, 0.820)
\end{equation}

The two endpoints are clearly distinct, confirming
Equation~\eqref{eq:noncommute}: identical primitives, identical
resultant $R$, opposite structural readings. Notably, if $\chi_{\text{Sun}}
= 0^\circ$ and $\chi_{\text{Wind}} = 90^\circ$, the accumulated phase
shift is $270^\circ$ in \emph{both} orderings, since phase addition is
commutative. The two poems therefore share the same blended color/mood
even as their spatial --- structural --- meaning diverges. This
asymmetry (shared affect, divergent structure) is a direct and
testable prediction of the model.

\subsection{Scope}
\label{subsec:transform_scope}
$T$ is deliberately left unconstrained beyond
Equation~\eqref{eq:transform_def}: any sequence of any transforming
states is a formally valid trajectory. Pattern Arithmetic in this form
supplies the substrate on which expression is built --- the
combinatorial space of possible trajectories --- and does not yet
supply a grammar constraining which trajectories are coherent,
well-formed, or meaningful to a human reader. Whether such constraints
belong inside this framework (as further axioms on admissible
sequences of $T$) or outside it (as a learned or externally specified
grammar layered on top) is left open.

\section{Analogical Arithmetic: The Offset Operator}
\label{sec:analogical_offset}

The $\oplus$ operator mints composite nouns, and the $T$ operator mints ordered sentences. However, analogical reasoning requires a third, distinct mathematical operation: the offset. By mapping the dictionary onto our framework such that each word is a labeled state on a single shared globe, relationships between concepts are captured by their relative spatial differences.

To assert that ``man is to king as woman is to queen,'' we do not average the vectors (which $\oplus$ handles), nor do we spin one around the other (which $T$ handles). Instead, we extract the relational offset between the first pair and apply it to the target:
\begin{equation}
    \text{king} - \text{man} + \text{woman} \approx \text{queen}
    \label{eq:analogical_offset}
\end{equation}
This operation reuses the geometric translation that turned man into king, demonstrating that the algebra handles analogical equivalence through vector differences. It is an analogical reuse of choreography, rather than a blending of paint.
To visualize this offset on a 2D planar cross-section of $\mathcal{S}$, let $\text{man} = 0^\circ$, $\text{woman} = 90^\circ$, and $\text{king} = 30^\circ$. The translational offset from $\text{man}$ to $\text{king}$ is $+30^\circ$. Applying this exact relational offset to $\text{woman}$ yields $90^\circ + 30^\circ = 120^\circ$, which defines the spatial coordinate for $\text{queen}$.
% TODO [T5] This example places man/woman/king/queen all on a single
% great circle, so "offset" reduces to ordinary real-number addition
% of one angle. That sidesteps the actual difficulty: subtraction and
% addition of points are not natively well-defined operations on a
% curved 2-sphere when the points are NOT coplanar/co-circular. Before
% claiming Eq.~\eqref{eq:analogical_offset} generally, define the
% offset construction for the general case -- e.g. log-map the two
% source points to a common tangent space, take the difference there,
% then exp-map that tangent vector at the target point -- and show a
% worked example where the three points do NOT share a great circle.

\section{Encodings, Sorts, and Mixed Documents}
\label{sec:encodings_sorts}

The unit probability sphere is a single bounded \emph{register}.
Numbers, words, and colors are not three different globes; they are
three encodings of a token onto the same globe. A scientific paragraph
already mixes those alphabets on one page. Pattern Arithmetic may do
the same, provided each loaded state carries a \emph{sort} and the
operators respect that sort.

\subsection{Two primary sorts}
\label{subsec:sorts}
A state $x\in\mathcal{S}$ is always a 4-tuple
$(r,\theta,\phi,\chi)$. What the tuple \emph{means} depends on the
encoding used to write it:
\begin{itemize}
    \item \textbf{Numeric sort ($\mathtt{num}$).}
          The state represents a scalar. One toy encoding stores an
          integer $n$ on a marked circle as
          $\chi = 2\pi n/N$ (arithmetic modulo $N$).
          Another, stereo-style, sends $x\to\pm\infty$ toward the
          poles so that unbounded reals still occupy the globe.
    \item \textbf{Lexical sort ($\mathtt{word}$).}
          The state is a labeled point in a toy dictionary living on
          the same globe. Similarity is the angle between points
          (cosine similarity). A bag of words is formed by $\oplus$.
          A sentence or poem is formed by $T$.
\end{itemize}
Schoolbook addition lives only in $\mathtt{num}$ mode:
\begin{equation}
    1 + 1 = 2, \qquad 7 + 5 = 12.
\end{equation}
Pattern blend lives only in $\mathtt{word}$ (or color) mode:
\begin{equation}
    A \oplus A = A.
\end{equation}
The two equations may share a sphere. They may not share an operator.
Pressing $\oplus$ on the encodings of $7$ and $5$ does not yield $12$;
pressing schoolbook $+$ on $\mathrm{sun}$ and $\mathrm{wind}$ does not
yield a chord.

\subsection{A document is a sequence, not a blend}
\label{subsec:document_sequence}
Words and numbers coexist in reports because a document is an ordered
tape of tokens, not one mixed arrow. Write a paragraph as a trajectory
of sorted states
\begin{equation}
    \pi = \bigl((x_1,\tau_1),\ldots,(x_m,\tau_m)\bigr),
    \qquad \tau_i\in\{\mathtt{num},\mathtt{word}\},
    \label{eq:mixed_document}
\end{equation}
exactly as a poem was a trajectory in
Section~\ref{subsec:transform_trajectories}. The sphere holds one
current occupant of the register; the paragraph is the history of
loads and actions.

They may meet in only three honest ways.

\paragraph{Juxtaposition.}
The default. Tokens sit in order and no algebra is applied between
them. In ``the sample reached $37.4^\circ\mathrm{C}$ at dawn,'' the
number $37.4$ is not blended with the word $\mathrm{dawn}$.

\paragraph{Typed action.}
The operator $T_w(x)$ is allowed across sorts only with an explicit
output sort. Two useful conventions, which the rest of this paper
adopts, are:
\begin{itemize}
    \item a word acting on a number yields a number-in-context
          (still $\mathtt{num}$);
    \item a number acting on a word yields a refined word
          (still $\mathtt{word}$).
\end{itemize}
A number may also set a numeric dial of $T$ --- for example the
rotation angle $\alpha(r)$ --- because that dial is already a scalar.
That is how ``wind of $12\,\mathrm{m/s}$'' can strengthen a turn
without becoming a synonym of $\mathrm{wind}$.

\paragraph{Cast.}
Rare and named. The map $E_{\mathtt{num}\to\mathtt{word}}(2)=\texttt{two}$
turns a digit into a spoken numeral; the inverse is parsing. Casts are
dictionary maps, not $\oplus$.

The illegal meeting is an untyped blend
\begin{equation}
    \mathrm{dawn} \oplus 37.4,
\end{equation}
which has no meaning until one side is first cast into the other's
sort.

\subsection{Operator legality}
\label{subsec:type_table}
Table~\ref{tab:type_table} records which buttons accept which pairs.
Here $+$ is schoolbook addition on the numeric encoding, $\oplus$ is
idempotent blend, $T$ is typed action, and $x-a+b$ is the analogical
offset of Section~\ref{sec:analogical_offset}.

\begin{table}[h]
\centering
\begin{tabular}{lcccc}
\hline
Pair & $+$ & $\oplus$ & $T_w(x)$ & $x-a+b$ \\
\hline
$(\mathtt{word},\mathtt{word})$ & no & yes & yes & yes \\
$(\mathtt{num},\mathtt{num})$   & yes & no  & no$^{*}$ & no$^{*}$ \\
$(\mathtt{word},\mathtt{num})$  & no & no  & typed & no \\
$(\mathtt{num},\mathtt{word})$  & no & no  & typed & no \\
\hline
\end{tabular}
\caption{Sort discipline on a single sphere.
$^{*}$Numeric $T$ and numeric offset may be defined later as
maps of the chosen numeric encoding; they are not the lexical
operators of Sections~\ref{sec:transformation_algebra}
and~\ref{sec:analogical_offset}.}
\label{tab:type_table}
\end{table}

\subsection{Worked mixed sentence}
\label{subsec:mixed_sentence}
Take the report line \emph{Wind of $12\,\mathrm{m/s}$ darkened the sun.}
The register is used as follows.

\begin{enumerate}
    \item Load $\mathrm{wind}$, sort $\mathtt{word}$.
    \item Load $12$, sort $\mathtt{num}$. Juxtapose; do not form
          $\mathrm{wind}\oplus 12$.
    \item Load $\mathrm{m/s}$ as a unit-name (lexical qualifier of the
          number).
    \item Load $\mathrm{sun}$, sort $\mathtt{word}$.
    \item Let the number cap a strength,
          $\sigma(12)\in[0,1]$, and set
          $r_{\mathrm{wind}}=\sigma(12)$ so that
          $\alpha=\pi\sigma(12)$ in Equation~\eqref{eq:transform_def}.
    \item Apply $T_{\mathrm{wind}}(\mathrm{sun})$. Output sort:
          $\mathtt{word}$.
\end{enumerate}
The scalar never became a hue. It turned a dial of an action whose
actors remained words. That is the same division of labor a printed
paper already uses: digits set magnitudes; words name what the
magnitudes belong to.

\subsection{Several lexicons, one concept sphere}
\label{subsec:shared_lexicons}
Surface tokens are not meanings. Two dictionaries may write the same
pair of concepts with unrelated strings:
\begin{equation}
    \begin{aligned}
        L_1&: \quad \texttt{xyz}\mapsto\mathrm{man},\quad
                    \texttt{qyz}\mapsto\mathrm{woman},\\
        L_2&: \quad \texttt{fgh}\mapsto\mathrm{man},\quad
                    \texttt{hgf}\mapsto\mathrm{woman}.
    \end{aligned}
    \label{eq:two_lexicons}
\end{equation}
In the usual pipeline each lexicon is tokenized, encoded, and
\emph{embedded separately}. The geometry that makes
$\mathrm{man}$ near $\mathrm{king}$ is then learned twice, and a
further alignment step is required before $L_1$ and $L_2$ can share
an analogy. The token-to-embedding map is doing two jobs at once:
naming a string, and inventing a place for its meaning.

The unit probability sphere is offered as a shared
\emph{concept register}, not as a replacement for a lexicon table.
A token is a surface form. A state on $\mathcal{S}$ is a meaning.
Each lexicon is only a cast
\begin{equation}
    E_{L}(\texttt{token}) = x\in\mathcal{S},
    \label{eq:lexicon_cast}
\end{equation}
the same kind of map as $E_{\mathtt{num}\to\mathtt{word}}$ above.
Under that discipline both $\texttt{xyz}$ and $\texttt{fgh}$ load the
\emph{same} state $\mathrm{man}$. The offset
$\mathrm{king}-\mathrm{man}$ of Section~\ref{sec:analogical_offset}
is then computed once, on the globe, and reused by every lexicon that
points at those two points.

Embeddings therefore stay. They do not disappear.
What changes is that they are no longer private vector spaces, one
per lexicon. They are embeddings \emph{into the same unit}
$\mathcal{S}$. Once $E_{L_1}(\texttt{xyz})=E_{L_2}(\texttt{fgh})=\mathrm{man}$,
the two surface forms occupy one address, and every angle measured
from that address is the same angle:
\begin{equation}
    \angle(\mathrm{man},\mathrm{woman})
    \;=\;
    \angle\bigl(E_{L_1}(\texttt{xyz}),\,E_{L_1}(\texttt{qyz})\bigr)
    \;=\;
    \angle\bigl(E_{L_2}(\texttt{fgh}),\,E_{L_2}(\texttt{hgf})\bigr).
    \label{eq:uniform_angles}
\end{equation}
The embeddings are uniform: same globe, same points for the same
concepts, same angles for the same relations. A raw string still
needs the lexicon map $E_L$. It no longer needs a new geometry.

\subsection{Deferred programme: standardized concept dictionaries}
\label{subsec:standard_dictionaries}
The present section only \emph{names} the shared globe. It does not
build the dictionaries. The later refinement, if the idea is pursued,
is a pair of catalogues rather than a new operator:
\begin{itemize}
    \item one \emph{concept dictionary} --- a stipulated inventory of
          unit states on $\mathcal{S}$ (or a higher-dimensional
          analogue), so that $\mathrm{man}$ and $\mathrm{woman}$ have
          standard addresses and therefore standard angles;
    \item one \emph{language lexicon} $E_L$ per written or spoken
          system, sending that language's tokens onto those addresses;
    \item learnable paths --- the species of $T$ and the dials inside
          a given $T$ (angle map, phase map, paint-versus-poem gate)
          may be fitted from examples. Learning would adjust the path,
          not the rule that a bowl is collected before it is read.
          See also Section~\ref{sec:conclusion}.
\end{itemize}
Standardization of the first catalogue is what makes embeddings
uniform across English, a cipher, or a second natural language.
The second catalogue is still a tokenizer-plus-encoding: it is the
phone book, not the city. Tokens and encodings have already earned
their keep, so they stay. They only have to work \emph{lighter}:
look up a standard address instead of inventing one. They disappear as
\emph{inventors of geometry}. They do not disappear as \emph{names}.
A world in which every language pointed at the same concept
dictionary would still print letters on a page; It would no longer need each 
corpus to learn a private $\mathrm{man}$--$\mathrm{woman}$ angle. 
That programme is recorded here and left open. It is not an
axiom, and it is not implied by the existence of $\mathcal{S}$ alone.

\subsection{Scope of the register model}
\label{subsec:register_scope_final}
This section claims only that one sphere can \emph{host} both alphabets
as a workspace, and that several lexicons can point at the same
workspace through different casts. It does not claim that the two-dimensional surface
$S^2$ is a realistic embedding space for a full natural-language
lexicon: a working word embedding typically needs tens or hundreds of
dimensions. The construction here is a typed register and a toy
dictionary on one globe. Scaling the lexical encoding is a later
choice --- a higher-dimensional sphere, a spherical-harmonic field, or
an external lookup table that writes states into the register. The
lookup tables $E_L$ may still be learned; the claim is only that they
should write into one geometry rather than each minting a private one.

\section{A Sketch of Conventions}
\label{sec:conventions_sketch}
These notes are a seed for later grammars, not a code to enforce.
They show one way a typed register \emph{might} be spoken. A later
grammar --- learned, cultural, or formal --- may replace every line.
This paper only needs the ground level enough that such a grammar has
somewhere to stand.

\paragraph{Case does not move the address.}
The tokens $\texttt{man}$, $\texttt{Man}$, and $\texttt{MAN}$ load the
same direction $(\theta,\phi)$. Orthography is a doorbell, not a new
city block.

\paragraph{Language does not move the address.}
English $\texttt{man}$ and a Polish token for the same concept load the
same point. Only the lexicon map $E_L$ changes. Language is a
phone-book tag $L$, not a second angle.

\paragraph{Color is reserved for mood, not for language.}
Phase $\chi$ is already used as affect in the poem example
(Section~\ref{subsec:transform_example}). If language or case also
occupied $\chi$, that example would break. Emphasis may tint $\chi$
or raise $r$; it should not relocate $\mathrm{man}$ to $\mathrm{king}$.

\paragraph{A bag is not a sentence.}
$\{\mathrm{sun},\mathrm{wind}\}$ is $\oplus$.
``Wind darkens the sun'' is $T$. Ingredients and order are different
verbs. How often a unit is used is not a second shell in the bag:
the path may visit $\texttt{twinkle}$ twice; $\mathtt{num}$ may
count $1+1+1+1=4$. $\oplus$ will not.

\paragraph{A numeral is not a word until it is cast.}
The digit $2$ stays $\mathtt{num}$. The word $\texttt{two}$ is
$\mathtt{word}$. There is no silent $\oplus$ between them
(Section~\ref{sec:encodings_sorts}).

A toy row makes the first three conventions visible. Place
$\mathrm{man}$ at $(\theta,\phi)=(10^\circ,40^\circ)$:

\begin{center}
\begin{tabular}{llcc}
Token & $L$ & $(\theta,\phi)$ & $\chi$ (mood only) \\
\hline
\texttt{man} & English & $(10^\circ,40^\circ)$ & $0^\circ$ (neutral) \\
\texttt{Man} & English & $(10^\circ,40^\circ)$ & $0^\circ$ \\
\texttt{MAN} & English & $(10^\circ,40^\circ)$ & $20^\circ$ (emphasis, optional) \\
\texttt{man} & Polish  & $(10^\circ,40^\circ)$ & $0^\circ$ \\
\end{tabular}
\end{center}

Same block. Different doorbells. Nothing here is binding.

\section{Operational Algebra \& Numerical Demonstration}

\subsection{Vector Addition \& Phase Blending}
Using Equation~\eqref{eq:interference_formal} from Axiom~\ref{subsec:axiom_interference}, combining Red Vector ($0^\circ$) and Blue Vector ($120^\circ$) at coordinate ($0.6, 0.8$) results in deterministic phase blending ($\chi = 66.6^\circ$).
% TODO [T4] Equation~\eqref{eq:interference_formal} produces a
% MAGNITUDE (I get 0.721 for these inputs), not a phase -- it cannot
% by itself yield "chi = 66.6 deg." Show the actual formula used to
% get 66.6 deg. I tried the obvious candidates -- linear-weighted
% average (68.6 deg), magnitude-squared-weighted average (43.2/76.8
% deg), and phasor/vector-sum angle (73.9 deg) -- none match. Either
% this number comes from an unstated rule that needs to be written
% down, or it's a placeholder that should be recomputed once T4 above
% is resolved.

\subsection{Rotational Transposition (Multiplication) \& Unmixing (Division)}
\begin{itemize}
    \item \textbf{Multiplication ($\times$):} State transformation via Lie group rotation matrix $\mathbf{R} \in \mathrm{SO}(3)$.
    \item \textbf{Division ($\div$):} Deconvolution and signal unmixing via inverse transformation $\mathbf{R}^{-1}$ or anti-phase isolation.
\end{itemize}
% TODO [T6] "Multiplication" here just restates Axiom 6 / the T
% operator under a new name -- no new content. "Division" asserts
% deconvolution/unmixing without a formula: given a combined state
% A+B (Axiom 5's magnitude-only output, itself incomplete per T4),
% what is the actual inverse operation that recovers A and B
% separately? Either give the formula or cut this until T4 is fixed,
% since division here depends on interference being fully defined.

\subsection{Computational Efficiency \& Spectral Compression}
\begin{itemize}
    \item Mitigating $O(N^3)$ voxel grid costs using Spherical Harmonics ($Y_l^m$) coefficient arrays $O(L^2)$.
\end{itemize}
% TODO [T6] One-line claim, no derivation: which quantity is being
% expanded in spherical harmonics (the field P from Axiom 1? a
% collection of point states?), what does L need to be to represent a
% single sharp state faithfully, and what's actually lost/gained vs.
% storing (r,theta,phi,chi) directly for one pattern? As written this
% reads as a citation to a known technique rather than a result of
% this paper's framework.

% TODO [T6] This entire section is single-bullet gestures with no
% connection back to the axioms of Sec. 2-4: no state assignment for
% "flora," "fauna," "soil," etc., no explanation of why ecosystem
% interaction should be $A+B=AB$ (undefined notation -- concatenation?
% multiplication in some encoding?), and no worked numeric example
% like Sec. 3's Sun/Wind/Sky case. As written these read as evocative
% analogies rather than applications of the formal machinery. Either
% develop one of these with an actual worked example (states,
% operators, a checkable numeric result), or move the section to a
% clearly-labeled "future applications" list until that's done.
\section{Real-World Physical \& Biological Models}
\subsection{Forest Ecosystem Dynamics}
\begin{itemize}
    \item Modeling flora, fauna, soil, and moisture resource completion via vector addition ($A + B = AB$).
    % TODO: "$A+B=AB$" is undefined notation -- is AB concatenation,
    % a new composite state, matrix product? Doesn't match any
    % operator defined in Sec. 2-4.
\end{itemize}

\subsection{Targeted Viral/Pathogen Cancellation (Selective Subtraction)}
\begin{itemize}
    \item Isolating disease signatures and applying $180^\circ$ anti-phase mirror vectors to achieve zero-state equilibrium ($A - B = 0$).
    % TODO: recall from T4 that A-B=0 in Axiom 5 requires |A|=|B|
    % exactly, not just opposite phase. A real pathogen signature and
    % a designed countermeasure won't generically have equal
    % magnitude -- say how magnitude-matching is achieved, or this
    % doesn't actually reach the null state.
\end{itemize}

\subsection{Atmospheric Storm Systems (Phase Shifts and Shear Turbulence)}
\begin{itemize}
    \item In-phase wave coherence vs. minor phase shift ($0.01$) driving localized vortices and trajectory shifts.
    % TODO: no mechanism connects a "phase shift" in this framework's
    % chi to vorticity/shear in an actual fluid-dynamics sense --
    % currently just a suggestive juxtaposition of vocabulary.
\end{itemize}

\section{Discussion \& Related Work}
% TODO [T7] All three subsections below are empty headers. At minimum
% this section needs actual comparisons to:
%  - RotatE (Sun et al.) -- knowledge-graph relations modeled as
%    rotations in complex/unit-circle space, structurally very close
%    to the T operator of Sec. 3. What's actually new here relative
%    to it?
%  - The Bloch sphere itself (named below but undiscussed) -- a qubit
%    is (theta,phi) plus a global phase on a unit sphere; explain the
%    genuine differences (e.g. this paper's r is not a Bloch-sphere
%    quantity, chi here is not a *global* phase) rather than just
%    naming the comparison.
%  - von Mises-Fisher / directional statistics on spheres, which
%    already formalizes "probability density over directions" --
%    relevant directly to the Axiom 1 field definition (see T1).
\subsection{Comparison to Quantum Bloch Spheres and Hilbert Spaces}
\subsection{Integration with machine learning latent embeddings (Cosine Similarity)}
\subsection{Hardware Realization: Optical Phase Computing vs. Digital GPUs}

\section{Conclusion \& Future Directions}
\label{sec:conclusion}

This paper has sketched a single pipeline: collapse many sorted units
into one state on a unit probability sphere, then derive results by
ordered transformation of that state. Embeddings, on this view, are
not abolished. They are embeddings into one geometry, so that a
concept such as $\mathrm{man}$ may keep the same angles under every
lexicon that points at it.

A few spoken conventions are sketched in
Section~\ref{sec:conventions_sketch} as seed, not statute.

What remains to be built, and is only filed here, is a standardized
concept dictionary together with per-language lexica $E_L$
(Section~\ref{subsec:standard_dictionaries}). Until those catalogues
exist, tokenization and encoding remain the handles by which a
paragraph finds its points. They may continue; they will only have to
work lighter. The open bet is that they need not also
remain the architects of a new space of meanings.

A further deferred question is whether a path $T$ should stay
hand-specified or carry learnable weights: the angle law
$\alpha(r)$, a phase map, a gate between $T_{\mathrm{paint}}$ and
$T_{\mathrm{poem}}$, and the lexicon casts $E_L$. The honest
constraint is structural. If undo-the-poem matters, learning should
stay inside $\mathrm{SO}(3)$, not collapse $T$ into an arbitrary
matrix. Weights may change how the shells are thrown. They should
not let the shells throw themselves.

Ambient $\mathbb{R}^3$ makes the expressions
$\mathcal{S}+(a,b,c)$ and $\mathcal{S}\cdot(a,b,c)$ look legal.
They are legal as \emph{vector} arithmetic. They are not automatically
legal as \emph{unit} arithmetic. A state with $r=0.98$ added to
another with $r=0.98$ along the same direction is $r=1.96$, which
has left the ball. A scalar $2\cdot(0.98)$ does the same. The unit
cap $r\in[0,1]$ then forces a choice, to be made later and not
here: clip (overflow saturates at $1$), project (renormalize
direction and forget extra strength), or refuse (the operator is
undefined off the ball). Schoolbook $+$ on the $\mathtt{num}$
encoding is a different door: decode, add in $\mathbb{R}$, encode
back if the chart allows. Do not let $\mathbb{R}^3$ addition
quietly replace $\oplus$, $\mathrm{Mix}$, or $T$. The embedding
lives in $\mathbb{R}^3$. The unit lives in the ball.

\end{document}